\section{Introduction}\label{sec:intro}
The APEIRON Framework~\cite{Beniers2026Apeiron} is a multi‑layer architecture for autonomous, non‑anthropocentric knowledge genesis.
Its foundational layer, Layer~1, defines \emph{Ultimate Observables} — elementary data atoms that are considered irreducible under any further decomposition.
In earlier work~\cite{Beniers2026Categorical}, we gave a categorical semantics for Layer~1: the Apeiron category \(\mathcal{A}\) has as objects the set \(\mathcal{O}\) of all observables and as morphisms the decomposition paths induced by four independent atomicity frameworks.
The central result of that paper is that an observable is multi‑axially atomic if and only if it is a \emph{separated object} in \(\mathcal{A}\) — there are no non‑identity morphisms out of it.

Layer~2 of the framework takes a collection of observables and organises them into a \emph{relational hypergraph}: a weighted undirected graph where vertices are observables and edges represent co‑activation strength.
Empirically, this hypergraph has been validated on images and time‑series data, spontaneously yielding functionally meaningful clusters.

In this paper we study a different structure that naturally arises from the decomposition operators: the \emph{decomposition hypergraph} of an observable.
This is the directed graph whose vertices are all observables that can be further decomposed from the given one, and whose edges are the decomposition steps.
We construct a category \(\mathcal{H}\) of finite directed graphs and a contravariant functor
\[
\mathbf{D} : \mathcal{A}^{\mathrm{op}} \longrightarrow \mathcal{H}
\]
that sends each observable to the subgraph induced by the vertices that are reachable from it via a directed decomposition path.
We prove:
\begin{enumerate}
\item If \(o\) is a separated object in \(\mathcal{A}\) (i.e.\ multi‑axially atomic), then \(\mathbf{D}(o)\) consists of a single isolated vertex.
\item For any finite set \(S \subset \mathcal{O}\), the union of the vertex sets of \(\{\mathbf{D}(o)\}_{o \in S}\) equals the set of all observables that can be reached from \(S\) by decomposition.
\end{enumerate}
The functor \(\mathbf{D}\) therefore faithfully translates the decomposition structure of Layer~1 into the language of graphs, providing a rigorous categorical lens on the “downstream decomposition” concept.
We verify the construction in Python using the reference implementation of Layer~1 and report the outcome for a small benchmark of observables.

The relation between \(\mathbf{D}\) and the relational hypergraph of Layer~2 is not an identity; rather, the decomposition graph informs a complementary aspect of the observable’s structural context.
A full categorical modelling of Layer~2 remains an open challenge.